\chapter{Referencial teórico}
\label{cap:referencial}

\section{O estatuto deste capítulo}
\label{sec:ref-estatuto}

A investigação não partiu de um arcabouço teórico em busca de confirmação
empírica. Partiu de uma pergunta sobre a dinâmica da pastagem em Goiás e
avançou pela curiosidade: cada resposta abria a pergunta seguinte, e foi essa
cadeia --- não uma hipótese teórica prévia --- que produziu as quatro frentes
da investigação e a tese da ``marcha ao norte''. A literatura foi mobilizada,
nesse percurso, sobretudo para decidir \emph{como perguntar}: qual estimador
resiste a séries integradas, qual teste de quebra não depende de datas
escolhidas pelo pesquisador, que unidade territorial permite comparar
municípios ao longo de quatro décadas. É desse uso que vem o instrumental do
Capítulo~\ref{cap:metodologia}.

O estágio atual é outro. Existem respostas, e falta a leitura sistemática que
diga o que a literatura já sabe sobre elas. Essa leitura é trabalho da
dissertação. A Seção~\ref{sec:ref-confrontados} reúne os arcabouços que os
dados já confrontaram --- cada um com uma previsão específica que a
investigação pôde testar, inclusive quando o teste a contrariou --- e a
Seção~\ref{sec:ref-agenda}, os que permanecem como agenda de leitura.

\section{A tese e os arcabouços que ela mobiliza}
\label{sec:ref-tese}

A tese que este trabalho defende pode ser enunciada em uma frase: Goiás
viveu, entre 1985 e 2024, uma reorganização espacial da produção
agropecuária --- intensificação no Sul, fronteira no Norte --- compatível com
a ação de forças de mercado comuns sobre um gradiente de aptidão, e limitada
pelo encolhimento da terra ainda exposta à conversão; e \emph{não} um
deslocamento causal de uma região sobre a outra. Das três cláusulas, a
segunda é a que o trabalho estabelece com mais firmeza --- é a que se testa
diretamente e cujo resultado é negativo ---, enquanto a leitura de mercado
comum permanece corroborante e não estabelecida.

Cada termo dessa frase pede um arcabouço próprio: \emph{fronteira} e
\emph{reorganização} remetem à teoria da fronteira agrícola; \emph{gradiente
de aptidão}, à economia da renda da terra; \emph{deslocamento causal}, à
literatura de mudança indireta de uso da terra (iLUC) e de teleacoplamento;
\emph{forças de mercado}, aos drivers macroeconômicos das commodities; e
\emph{teto de oferta}, à teoria da transição florestal. Nem todos foram
percorridos com a mesma profundidade, e a seção seguinte trata apenas
daqueles em que houve confronto efetivo com o dado.

\section{Os arcabouços e as previsões que eles geram}
\label{sec:ref-confrontados}

O critério de inclusão nesta seção é estreito e deliberado: entra o arcabouço
do qual foi possível extrair uma previsão específica, verificável na série de
Goiás. O que se enuncia aqui é a previsão; o resultado de cada confronto ---
que em dois dos cinco casos não foi favorável ao arcabouço --- está no
Capítulo~\ref{cap:discussao}, e cada subseção indica onde.

\subsection{Renda da terra e localização}
\label{sec:ref-renda}

A base espacial do argumento é o modelo de localização agrícola de von Thünen
(1826), para o qual a terra é alocada ao uso que lhe rende mais, e a renda de
cada uso é determinada pela localização. Dessa formulação decorre uma
previsão testável: a expansão agrícola ocupa primeiro as terras de melhor
aptidão e só depois as piores --- e, portanto, os usos devem ordenar-se no
espaço segundo o gradiente de aptidão.

Testar isso exige medir o gradiente por fora do fenômeno que ele deveria
explicar. Inferir aptidão a partir do uso observado seria circular; a
alternativa adotada, um mapa edafoclimático construído sobre atributos
físicos, é descrita no Capítulo~\ref{cap:metodologia}.

\citeonline{Angelsen2007} oferece a síntese que amarra a dimensão espacial à
temporal, ao combinar explicitamente von Thünen com a transição florestal: a
renda agrícola crescente eleva o desmatamento até que o próprio
desenvolvimento a reduza, ou até que a escassez de vegetação eleve a renda do
remanescente e reverta a trajetória. A previsão que daí se extrai é a de que
o gradiente espacial e o pico temporal sejam faces do mesmo processo de renda
da terra, e não dois fenômenos distintos. O confronto está na
Seção~\ref{sec:disc-fronteira}.

\subsection{Transição florestal}
\label{sec:ref-transicao}

A teoria da transição florestal \cite{Mather1992, Rudel2005} descreve uma
trajetória em três fases: cobertura alta e desmatamento baixo, conversão
acelerada, e por fim estabilização ou ganho líquido de vegetação. É o
arcabouço com que a pesquisa lê o ``teto de oferta'': o Sul do estado, que
abriu cedo, deveria comportar-se como região em fase tardia, enquanto no
Norte resta Cerrado ainda exposto à conversão.

A previsão que interessa é a da fase tardia, e ela tem duas metades: estoque
reduzido com taxa de conversão cadente, \emph{e} uma virada de sinal ---
recuperação, regeneração líquida da vegetação. As duas metades são
separáveis no dado, e essa separação é o que a Seção~\ref{sec:disc-fronteira}
examina. Se a resposta for peculiaridade do Cerrado, do período ou do estado,
este desenho não decide --- pergunta que a agenda da
Seção~\ref{sec:ref-agenda} precisa endereçar.

\subsection{Deslocamento indireto e teleacoplamento}
\label{sec:ref-iluc}

A literatura de mudança indireta de uso da terra estabelece que a expansão da
soja desloca a pecuária e, com ela, a fronteira. \citeonline{Lapola2010}
mostram que esse efeito indireto pode anular a economia de carbono dos
biocombustíveis brasileiros, ao empurrar a fronteira de pastagem para as
florestas amazônicas e para o Cerrado; \citeonline{Arima2011} oferecem
confirmação estatística do mecanismo na Amazônia, estimando que uma redução
de 10\% da soja em áreas de pasto consolidado teria reduzido o desmatamento
em até 40\% nos municípios mais florestados. O fenômeno é real na escala em
que foi estabelecido --- entre regiões e entre biomas.

A previsão que este trabalho pôde testar é mais estreita, e essa literatura
não a isola: o deslocamento \emph{intra-estadual} --- a lavoura do Sul goiano
empurrando o pasto para o Norte do próprio estado, por adjacência e com
precedência temporal curta.

Havia razão teórica para esperá-lo, e convém registrá-la antes do resultado.
\citeonline{Meyfroidt2010} definem \emph{leakage} como ``o deslocamento do
desmatamento para localidades vizinhas, por migração dos agentes do
desmatamento ou por comércio'': a vizinhança está no enunciado, e a migração
de agentes opera a curta distância. O arcabouço do teleacoplamento
\cite{Liu2013}, por sua vez, define seu objeto como as interações \emph{a
distância} e põe os acoplamentos locais de lado --- não prevê, portanto, a
ausência de efeito de vizinhança. E o precedente mais próximo,
\citeonline{Arima2011}, opera justamente no nível do município brasileiro,
com matrizes de peso espacial, e encontrou o efeito. O confronto está na
Seção~\ref{sec:disc-iluc}.

\subsection{Intensificação e disputa por terra}
\label{sec:ref-intensificacao}

Para a pecuária brasileira, \citeonline{Cohn2014} examinam, com um modelo
econômico de uso global da terra, se políticas de intensificação poderiam
poupar floresta e reduzir emissões. Dois de seus resultados oferecem previsão
verificável aqui. O primeiro é que a terra poupada fica bem aquém do
proporcional ao ganho de produtividade: a poupança existe, mas é parcial. O
segundo é locacional --- sob as políticas simuladas, a pastagem intensiva tem
cerca do dobro da chance da pastagem convencional de se instalar justamente
nas terras de melhor aptidão e logística para soja, de modo que a
intensificação da pecuária tende a \emph{competir} com a expansão da lavoura
pela mesma terra. Essa disputa, projetada ali em cenário, pode aqui ser
observada como desfecho medido no dado municipal
(Seção~\ref{sec:disc-iluc}).

\subsection{O câmbio e a divisão de trabalho com a literatura}
\label{sec:ref-cambio}

\citeonline{Richards2012} constituem a referência-âncora do driver cambial: a
desvalorização das moedas locais no fim dos anos 1990 elevou o retorno
esperado da soja e disparou resposta de oferta em área. Os autores estimam
que da ordem de 80~mil~km$^2$ --- 31\% da extensão então cultivada em Brasil,
Bolívia e Paraguai --- surgiram como resposta às desvalorizações do período;
só no Brasil, 63~mil~km$^2$, ou 29\% da área nacional de 2009.

O que este trabalho pode acrescentar a isso é limitado por construção, e o
limite precisa ser dito aqui para que o Capítulo~\ref{cap:resultados} seja
lido na chave certa. Com efeito fixo de ano, o painel absorve o nível do
choque comum e só deixa estimável o \emph{gradiente}: se o mesmo choque
desloca a atividade para onde a aptidão é \emph{menor}, com o núcleo apto
respondendo em lavoura e a fronteira, em gado. A divisão de trabalho é,
portanto, explícita: é a literatura que carrega o nível causal
câmbio$\to$soja$\to$fronteira, já ancorado em escala nacional e continental;
a tese contribui com a geografia desse efeito dentro do estado, e não
poderia contribuir com mais. O peso do que se obteve está na
Seção~\ref{sec:disc-drive}.

\subsection{O espelho do Cerrado}
\label{sec:ref-matopiba}

A literatura empírica mais diretamente comparável é a da fronteira do
Cerrado. \citeonline{Soterroni2019}, ao avaliar a extensão da Moratória da
Soja ao bioma, consolidam o Cerrado como fronteira ativa da agricultura
brasileira, com a expansão concentrando-se no norte e nordeste --- a região
do MATOPIBA (Maranhão, Tocantins, Piauí e Bahia), que nas projeções dos
autores absorve 86\% da expansão da soja no bioma até 2050 e concentra a
maior parte da conversão de vegetação natural para soja, por ser onde restam
os maiores remanescentes não perturbados. Como magnitude observada do
fenômeno, \citeonline{Rausch2019} estimam que a conversão direta para soja
respondeu por 16 a 32\% do desmatamento anual do Cerrado entre 2003 e 2014.

O achado-espelho decisivo para o posicionamento é a composição dessa
expansão, e ele está em \citeonline{CarneiroFilho2016}: de 2000 a 2014 a área
agrícola do Cerrado cresceu 87\%, e essa expansão se deu essencialmente sobre
áreas já antropizadas --- em torno de 70\% do total no período, chegando a
74\% na década final, com mais da metade sobre pastagem.
A exceção é o MATOPIBA que, ``por não dispor de áreas antropizadas com
aptidão para a agricultura, teve a maioria da expansão sobre áreas de
vegetação nativa''. Como a soja responde por 90\% da área agrícola do bioma,
a dinâmica descrita é, na prática, a da soja.

A previsão, para um estado que atravessa o mesmo bioma no eixo Sul--Norte, é
a de que esse contraste entre Cerrado maduro e fronteira apareça dentro de
Goiás --- e, aqui, sob um único desenho analítico, na mesma malha e com o
mesmo classificador (Seção~\ref{sec:disc-matopiba}).

\section{A linhagem do centróide migratório}
\label{sec:ref-centroide}

A métrica-manchete da primeira frente --- o centro de massa da conversão e
seu deslocamento --- pertence a uma linhagem metodológica consolidada, e não
é construção ad hoc. A elipse de desvio-padrão foi formalizada por
\citeonline{Lefever1926} e atualizada como instrumento de descrição espacial
por \citeonline{Yuill1971}; o centro mediano, obtido pelo algoritmo de
Weiszfeld, completa o instrumental clássico da estatística espacial
descritiva. Essa linhagem responde à objeção de ``métrica caseira'' e situa a
escolha metodológica do Capítulo~\ref{cap:metodologia}.

\section{Agenda de leitura para a dissertação}
\label{sec:ref-agenda}

Cada item traz a fonte identificada e a pergunta que a leitura deve decidir.

\begin{description}\itemsep2pt
\item[Intensificação induzida pela demanda.] Boserup (1965); Strassburg \emph{et al.}
(2014). \textbf{Decide:} se o regime do Sul é intensificação induzida por demanda ou
substituição por renda da terra, e se os dois são distinguíveis com os dados disponíveis.
\item[Deslocamento em escala subnacional.] Levantamento a partir de
\citeonline{Arima2011}, precedente já conhecido; Seto \emph{et al.} (2012) na fila.
\textbf{Decide:} se o nulo obtido aqui diverge de um resultado consolidado, replica um
nulo já reportado, ou não encontra par.
\item[Transição florestal, revisitada.] Literatura posterior a \citeonline{Rudel2005},
sobre savanas tropicais. \textbf{Decide:} se a estabilização sem recuperação observada no
Sul é caso conhecido do arcabouço ou observação a registrar.
\item[Fronteira na tradição brasileira.] Martins (1971); Becker. \textbf{Decide:} se as
categorias sociológicas admitem correspondência com cobertura do solo, ou se o uso é
analógico --- como a Seção~\ref{sec:disc-fronteira} já supõe.
\item[Composição da expansão no Cerrado.] Spera \emph{et al.} (2016); Noojipady \emph{et
al.} (2017). \textbf{Decide:} se as estimativas de \citeonline{CarneiroFilho2016}
convergem com as de bases independentes.
\item[Centróides migratórios em uso da terra.] \textbf{Decide:} se a métrica tem
precedente em matrizes de transição, o que mudaria como defendê-la.
\item[Câmbio, período recente.] Estudos do USDA-ERS sobre a depreciação de 2015--2016.
\textbf{Decide:} se o mecanismo de \citeonline{Richards2012} se repete no episódio
recente.
\item[Fronteira agrícola em Goiás.] Produção da UFG e da literatura regional.
\textbf{Decide:} o que já está estabelecido para Goiás, e portanto o que aqui é
contribuição.
\item[Von Thünen no original.] \textbf{Decide:} se a leitura corrente do modelo é fiel à
formulação.
\end{description}

\section{Síntese e posicionamento provisório}
\label{sec:ref-sintese}

O Quadro~\ref{quadro:posicionamento} resume o posicionamento na literatura
efetivamente conferida. É provisório por construção: a agenda da
Seção~\ref{sec:ref-agenda} pode alterar tanto as âncoras quanto a avaliação
do que a pesquisa acrescenta.

\begin{quadro}[htb]
\centering
\caption[Posicionamento provisório na literatura, por frente]{Posicionamento provisório na literatura conferida, por frente da
investigação.}
\label{quadro:posicionamento}
\small
\begin{tabular}{p{2.4cm}p{3.2cm}p{2.9cm}p{5.4cm}}
\toprule
\textbf{Frente} & \textbf{Arcabouço} & \textbf{Âncoras conferidas} & \textbf{O que o confronto mostrou} \\
\midrule
1. O padrão (a marcha) & centro de gravidade; renda da terra &
\citeonline{Lefever1926}; \citeonline{Yuill1971}; \citeonline{Angelsen2007} &
quantifica a marcha e o gradiente latitudinal persistente, com controle do
efeito de escala; a premissa de aptidão é medida por fonte exógena \\
\addlinespace
2. Mecanismo local duplo & intensificação $\times$ extensificação &
\citeonline{Cohn2014} &
a disputa por terra projetada em cenário aparece resolvida no dado municipal,
a favor da lavoura \\
\addlinespace
3a. Deslocamento (nulo) & iLUC; teleacoplamento; \emph{leakage} &
\citeonline{Lapola2010}; \citeonline{Arima2011}; \citeonline{Liu2013};
\citeonline{Meyfroidt2010} &
testa o canal em Goiás sob adjacência dirigida e não o encontra; o arcabouço
não previa esse nulo, e o precedente subnacional mais próximo
(\citeonline{Arima2011}) encontrou o efeito em outra região, desfecho e
especificação \\
\addlinespace
3b. Drive comum & drivers macro de commodities &
\citeonline{Richards2012} &
a literatura carrega o nível causal; o painel só estima o gradiente, e como
corroborante, não estabelecido \\
\addlinespace
4. Teto de oferta & transição florestal &
\citeonline{Mather1992}; \citeonline{Rudel2005} &
a previsão de recuperação na fase tardia \emph{não} se cumpre: estabilização
por depleção, sem regeneração líquida \\
\bottomrule
\end{tabular}
\fontefig{elaboração própria, a partir das obras conferidas em \texttt{ref/pdf/}.}
\end{quadro}

À pergunta ``o que há de novo, se a fronteira do Cerrado rumo ao norte já é
conhecida?'', a resposta provisória é: a literatura conferida documenta a
marcha da fronteira do bioma e o deslocamento indireto entre regiões; esta
pesquisa leva essas leituras à alta resolução intra-estadual, sobre malha de
território constante, e testa formalmente o canal de deslocamento local ---
que não se confirma ---, reconstruindo a interpretação como reorganização
coordenada por um drive macroeconômico comum sobre um gradiente de aptidão,
contra um teto físico de oferta de terra. Se essa resposta sobrevive ao
confronto com a literatura ainda não lida é, justamente, o que a agenda da
Seção~\ref{sec:ref-agenda} vai determinar.
